\section{Conclusion and Future Work}

This paper presented RETROSPECT, a Kubernetes-native secure Cloud--Fog--Edge platform for embedded WebAssembly life-cycle management. The platform combines CRD-based orchestration, gateway-centric TLS/CBOR communication, Zephyr-based device integration, and WAMR execution into a coherent architecture that bridges declarative cloud intent with constrained edge execution.

The most significant practical result is that the full workflow studied in this paper now completes reliably at the system level. In the final 30-trial campaign, enrollment, heartbeat supervision, deployment, and end-to-end transactional completion all achieved a success rate of 1.0, with sub-second mean latencies for enrollment and deployment and a mean end-to-end completion time of approximately 1~s. This finding matters because it confirms that reliable embedded orchestration depends not only on the correctness of the communication protocol, but also on the correctness of the policy used to interpret runtime evidence in the control plane. By documenting both the initial failure modes and the corrective strategy, the paper provides a more credible account of system behavior than a nominal-path-only evaluation would offer.

Taken together, the architecture, implementation, and validation show that cloud-native operational patterns can be adapted to embedded orchestration without discarding the security, observability, and reproducibility properties that make them valuable. The addition of transaction-oriented metrics, percentile-based latency analysis, and figure-reproducible experimental artifacts further strengthens the platform's value as a research baseline. RETROSPECT therefore offers a usable experimental baseline for future work on distributed cyber-physical software infrastructures in the ING-INF/05 domain.

Future work will extend the present platform in three directions. First, multi-gateway experiments with fault injection will be used to study placement, failover, and coordination at larger scale. Second, the northbound status path will be hardened so that CRD state and runtime state converge without ambiguity even when all controllers are active. Third, the evaluation will be expanded toward physical hardware, energy-aware measurements, and larger-scale workload mixes to complement the present system-level timing and transactional results.
